\section{Background \& Summary}

Land surface temperature (LST) describes the radiative thermal state of the land surface and is a key variable for understanding surface energy exchange and the water cycle \parencite{norman1995terminology,li2013lstReview,li2023globalLstReview}. Solar radiation, water availability and land cover jointly influence the partitioning of sensible and latent heat, producing thermal contrasts among landscapes \parencite{kustas2009thermal}. Leaf emergence, growth and senescence further alter canopy structure and its exchanges of heat and water with the atmosphere, giving surface thermal conditions a pronounced seasonal character \parencite{piao2019phenology}. Long-term satellite observations capture these contrasts from regional to continental scales and provide a spatial basis for studying the relationship between land cover and surface thermal environments.

During the leaf-on season, canopy development and water availability jointly influence evaporative cooling, linking LST to vegetation growth and water status \parencite{kustas2009thermal}. Temperature contrasts between cropland and surrounding surfaces therefore provide useful information on irrigation and crop-growing environments \parencite{anderson2012landsatThermal}, while contrasts between urban vegetation and built-up areas have also attracted substantial attention \parencite{voogt2003urban,weng2004lst}. Annual LST records focused on this season can support studies of cropland and oasis thermal environments, temperature differences around urban green spaces, and the relationship between vegetation cover and regional thermal patterns. The timing of vegetation growth varies among climatic regions \parencite{piao2019phenology}; a fixed seasonal window provides a common date range for comparing spatial thermal patterns across years.

Existing LST products provide extensive observations, but long-term coverage, temporal sampling and spatial detail remain difficult to achieve together. Kilometre-scale MODIS products frequently observe the regional thermal background \parencite{wan2008modis,duan2019modisValidation}, yet individual pixels often contain cropland, buildings, vegetation and bare ground. The Landsat archive spans several decades of finer-resolution thermal infrared observations, with native thermal-band resolutions of 120~m for Landsat 5, 60~m for Landsat 7 and 100~m for Landsat 8/9 \parencite{wulder2022landsat50}. Collection 2 surface temperature and 30 m gridded products for mainland China have advanced fine-scale thermal mapping over large areas \parencite{malakar2018landsatSt,usgs2021collection2,cheng2021landsat30m}. However, cloud cover, limited revisit frequency and intersensor differences lead to uneven annual distributions of valid observations, requiring reconstruction and compositing for long-term records \parencite{fu2016consistent}. Thermal sharpening and spatiotemporal fusion use finer-resolution auxiliary imagery to supplement spatial information \parencite{agam2007sharpening,qi2026thermalGap}. Applying these approaches to a long-term national annual record also requires continuous fine-scale inputs and clearly defined seasonal selection and compositing procedures.

This study combines long-term Terra MODIS thermal observations with 30 m annual leaf-on Landsat surface reflectance composites from the China Earth Observation Data Cube \parencite{cai2025ceodc}. The latter provide six bands---blue, green, red, near-infrared and two shortwave-infrared bands---from which we calculate optical indices. These indices, together with topographic and geographic variables, describe annual spatial differences in surface conditions. Annual thermal targets are constructed as the pixelwise median of quality-screened MODIS eight-day daytime composites whose nominal start dates fall within day of year (DOY) 150--300. The optical inputs are leaf-on composites, whereas the LST temporal window is explicitly defined by these dates. Random forest regression relates the annual thermal targets to surface predictors matched to the target spatial support \parencite{hutengs2016rf,ebrahimy2021adaptive}. A single cross-year model produces a 30 m temperature baseline, and smoothed kilometre-scale residuals then adjust the regional thermal background, combining annual thermal observations with local surface information.

We release \chinatherm{}, a record of annual 30 m leaf-on daytime LST composites across China for 2000--2024, comprising 25 annual products with 91 tiles of 4° by 4° per year. A common model and compositing procedure provide a long-term spatial record for studying the seasonal thermal environments of vegetation, cropland and urban landscapes. Blocked validation, cross-scale consistency analysis, and comparisons with Landsat 8, unmanned aerial vehicle (UAV) imagery and ground observations predominantly from the Heihe basin evaluate model generalization, national temperature patterns and local thermal conditions. Cases from the new urban districts of Yan'an and Lanzhou illustrate long-term thermal changes accompanying landscape development.

\begin{figure}[H]
\centering
\figureorplaceholder{../figures/main/Figure1.pdf}{}{7cm}
\ReviewFloatBegin
\caption{Study area, tiling scheme and observational support for the thermal target. a, Topography of China and the 91 output tiles of 4° by 4°. b, Multiyear median of the annual number of valid daytime eight-day thermal composites at each pixel during 2000--2024. c, Annual numbers of valid eight-day thermal composites and the proportion of land area with low observational support across China.}
\label{fig:study-area}
\ReviewFloatEnd
\end{figure}
